\paragraph{A theta coefficient attaining the valuation bound}
Write \(v_\ell\) for the valuation at \(\iota\) normalized by \(v_\ell(\ell)=1\).
The degrees \(l_\tau=c(\tau)-u\) satisfy
\(\kappa(\tau)\le l_\tau\le\kappa(\tau)+x^{42}\). Write \(l!=\prod_\tau l_\tau!\).
We first produce a coefficient of a theta kernel of rank \(12\) with valuation at
most \(F(\ell)+v_\ell(t)\). The coefficient includes \(l!\) and the integral
vacuum normalization. The scalar \(t\ne0\) is independent of the finite place
and all its conjugates have absolute value at most \(\exp(fmb/k)\).

Let \(\alpha_{\rm W}\) be the idele in the Weil representation and
\(\mathfrak a_0=E\cap\alpha_{\rm W}\widehat{\mathcal O}_E\) the associated lattice.
Take vectors \(w_i\) in \(\mathfrak a_0\) with multiplicities \(c_i\ge6fb\)
so that \(\sum_i c_i\tau(w_i)\overline{\sigma(w_i)}=M\delta_{\tau\sigma}\)
at all embeddings \(\tau,\sigma\).
We use the \leansrc{New/Derivatives/Vacuum/ShortVectors}{construction of the frame in the lattice},
which also gives, at each \(\ell\), a linear form whose coefficients are all units.
Denote by \(\mu_i\) the linear form with coefficients the conjugate embeddings
of \(w_i\) multiplied by the period normalization, and put
\(\Theta=\prod_i\theta(\mu_i z)^{c_i}\). Let \(H\) range over integral series in
the \(y(\mu_i z)\) of Section~\ref{sec:upper-bound}.
Write \(V_j(\nu)\le r\) to mean that in a product
\(\prod_{c=1}^j(H_c\Theta)(z_c)\) in \(j\) independent copies of the variables,
one can choose a coefficient of valuation at most \(r\) with combined multidegree \(\nu\).

First choose six \(\delta_c\in\mathbb Z_\ell\) with
\(\delta_0=1\) and \(\sum_c\delta_c^2=0\). Extend the four-square theorem to
\(\mathbb Z_\ell\) by compactness and write \(-1\) as a sum of four squares.
These four entries together with \(1,0\) give the required six.
With \(g=-\log\theta\), the expression \(\exp(\delta^2g(z)-g(\delta z))\)
is an integral series in \(y(z)\). This follows from the quotient identity
for the theta function under integer multiplication by taking coefficientwise
limits for \(\delta\in\mathbb Z_\ell\). Thus \(\prod_c\Theta(\delta_c z)\)
is a unit in the ring of \(H\)'s. Its constant term is \(1\), so its inverse
expands in the same ring.
\leansrc{New/Derivatives/Sharpness/LeastValuation}{The unit from the six theta factors}.

At every order \(N\), there is an integral series in \(y\) with
\(N\)-th derivative at zero of absolute value at least \(\ell^{-\gamma_\ell N}\).
Here \(\gamma_\ell=1/(\ell-1)-\rho_\ell\) at nonsplit primes and
\(\gamma_\ell=0\) at split primes. Indeed, let \(D_N\) be the largest absolute
value of the \(N\)-th derivatives of the \(y^j\). Since \(y(0)=0\), only
\(j\le N\) need be considered. The formal group law writes \(y(z_1+z_2)^j\)
as an integral series in \(y(z_1),y(z_2)\). Taking \(m_1,m_2\) derivatives
in the respective variables gives \(D_{m_1+m_2}\le D_{m_1}D_{m_2}\).
Let \(q_\ell\) be the norm of a prime of \(K\) above \(\ell\).
At \(N=q_\ell^e\), the \(N\)-th coefficient of the scaled formal parameter
is \(1\) in the residue field. Section~\ref{sec:upper-bound} gives the scaling
for nonsplit primes. For a split prime the same congruence holds with
\(q_\ell=\ell\), \(\rho_\ell=1/(\ell-1)\).
\leansrc{Lib/CM/FormalGroup/FormalParameter}{The formal parameter congruence at every prime}.
Undoing the scaling and multiplying by \(N!\), with
\(v_\ell(N!)\le(N-1)/(\ell-1)\), gives \(D_N\ge\ell^{-\gamma_\ell(N-1)}\).
Repeated first differentiation using the formal group law gives \(D_j\le1\).
For a fixed \(n\), write \(N=n\lfloor N/n\rfloor+(N\bmod n)\) and let
\(N=q_\ell^e\) tend to infinity. Submultiplicativity then yields
\(D_n\ge\ell^{-\gamma_\ell n}\).
\leansrc{New/Derivatives/Sharpness/OneCurve}{Derivative bound in one variable}.

Substitute the form \(\mu_i\) with unit coefficients in a series realizing
the bound for order \(|\nu|\). The multinomial expansion gives a coefficient
at \(\nu\) of valuation at most \(\gamma_\ell|\nu|-v_\ell(\nu!)\).
Multiply the first of the six factors by this series and by
\((\prod_c\Theta(\delta_c z))^{-1}\). After \(z_c=\delta_c z\), the product
equals the chosen series. The substitution is integral, so there is a
coefficient of at least this absolute value before substitution. Hence
\(V_6(\nu)\le\gamma_\ell|\nu|-v_\ell(\nu!)\).
At an inert prime the frame divided by \(\bar\pi_M\) resolves the identity.
It expresses a unit times each \(z_\tau\) as an integral linear combination
of the \(\mu_i(z)\). Replace the \(\mu_i(z)\) by \(y(\mu_i z)\). Each resulting
series has lowest term a unit times \(z_\tau\), so products give unit
coefficients in arbitrary multidegrees. With \(\Theta(0)=1\), this gives
\(V_6(\nu)\le0\) at inert primes.

\paragraph{Selection of couplings and factorials}
At a nonsplit prime write \(\tau_e=(-\ell^e)\tau\), with \(I_\ell\) as in
Section~\ref{sec:upper-bound}. First consider an inert prime. Select couplings
\(z_\tau^{\ell^{e_\tau}}z_{\tau_{e_\tau}}\) with \(e_\tau\) odd, using each
\(j_\tau\) times. Let \(\nu',\nu''\) be the multidegrees of the first and second
factors. In the choices below, a coupling from a selected starting place to
a selected endpoint with exponent \(e\le e_\tau\) must be the chosen coupling
for that endpoint. Use the normal form \(\Theta(\varpi z)=W\exp(C\Xi)\)
from Section~\ref{sec:upper-bound}, where \(C=\varpi^{\ell+1}/\ell\) has
\(v_\ell(C)=1/(\ell-1)-1<0\). Insert \(z'+z''\) in the first factor and
\(\delta_c z'\) in the other five and call the product \(G\).
At \(z''=0\) it is the integral unit above. Differentiate \(|j|=\sum j_\tau\)
times in the endpoint variables of \(z''\) and divide by \(C^{|j|}\).
Modulo the maximal ideal the result at \(z''=0\) is
\(\prod_\tau(\partial_{\tau_{e_\tau}}\Xi)^{j_\tau}\) times the unit.
Each differentiation on the exponential contributes \(C\); the terms
with fewer such factors have a factor \(1/C\) after division and vanish
in the residue field. In characteristic \(\ell\) the derivative of an
\(\ell\)-th power vanishes. Setting the \(z'\) at unselected starts to zero,
the lowest term of each \(\partial_{\tau_{e_\tau}}\Xi\) is a unit times
\(z_\tau'{}^{\ell^{e_\tau}}\). The coefficient at \((z')^{\nu'}\)
is therefore nonzero after reduction. Dividing by \(\nu''!\) and
undoing the scaling by \(\varpi\) gives
\begin{equation}\label{eq:attain-coupling}
 V_6(\nu'+\nu'')
 \le \left(\frac1{\ell-1}-1\right)|j|
      -\rho_\ell|\nu'+\nu''|-v_\ell(\nu''!).
\end{equation}
The right side is at most \(-|j|\) since
\((\ell+1)|j|\le|\nu'+\nu''|\). Also
\(v_\ell(\nu'!)\le|\nu'|/(\ell-1)\) and \(|\nu''|=|j|\), so
\(v_\ell(\nu'!\nu''!)+V_6(\nu'+\nu'')\le\gamma_\ell|\nu'+\nu''|-|j|\).
\leansrc{New/Derivatives/Sharpness/Couplings}{Differentiation at endpoints and the lowest term}.

Put \(L=\lceil6b/k\rceil\), \(l_\tau^-=\max(l_\tau-L,0)\), and
\(B_\tau=(l_\tau^-+1)\cdots(l_\tau^-+L)\).
Then \(L!\mid B_\tau\) and
\(v_\ell(l_\tau!)\le v_\ell(l_\tau^-!)+v_\ell(B_\tau)\).
The bound we need is
\begin{equation}\label{eq:attain-twelve}
 v_\ell(l!)+V_{12}(l)\le F(\ell)+2\sum_\tau v_\ell(B_\tau).
\end{equation}
At split primes \(F(\ell)=0\), and the bound follows from
\(V_6(l)\le-v_\ell(l!)\) by adjoining six factors of degree \(0\).

A nonsplit \(\ell>k\) is inert since \(\ell\ne d\).
Use \(e_\tau=1\) for all \(\tau\), and set
\(i_\tau=\lfloor l_\tau^-/\ell\rfloor\) and \(j_\tau=\min(i_\tau,l_{\tau_1})\).
The map \(\tau\mapsto\tau_1\) is bijective. At each coordinate the
starting degree is at most \(l_\tau^-\) and the incoming degree is at most
\(\min(L,l_\tau)\); their sum is at most \(l_\tau\).
Use the first six factors for the couplings and the remaining six for
the remaining degree at valuation at most \(0\). Then \(V_{12}(l)\le-\sum j_\tau\).
Since \(i_\tau\le L\), the factorial recursion gives
\(v_\ell(l_\tau!)\le i_\tau+2v_\ell(B_\tau)\).
Moreover \(L\ge x^{42}\), so \(i_\tau\le\kappa(\tau)/\ell\) and
\[
 i_\tau-j_\tau\le
   \bigl(\kappa(\tau)/\ell-\kappa(\tau_1)\bigr)_+ .
\]
The right side is the term of \(F(\ell)\) at \(\tau,e=1\).

At nonsplit \(\ell\le k\), use pairs with \(e\in I_\ell\), \(2e<l_0\), and both
\(\tau,\tau_e\) in \(\Phi_b\). If \(\ell=d\), then \(d\equiv q_0/2-1\pmod{q_0}\) and
\(d^2\equiv1\pmod{q_0}\). Thus \(-d^e\) is never \(1\) modulo \(q_0\),
and there are no such pairs. In the inert case sparsity makes the endpoints and starting places of
distinct pairs disjoint. Two outgoing pairs would give exponents \(0,e,e'\)
in the forbidden relation. Two incoming pairs give \(0,e-e',e\)
from the start of the pair with larger exponent \(e\).
For the end of one pair that starts another, use \(0,e,e+e'\).
A nontrivial return to the same place is ruled out by the order bound in
Section~\ref{sec:data}. We can therefore take
\(j_\tau=\lfloor l_\tau/\ell^{e_\tau}\rfloor\).
The starting degrees are at most \(l_\tau\), and the incoming degree
is at most \(2b/2\le l_{\tau_e}\). Put \(\eta=l-\nu'-\nu''\).
Apply \eqref{eq:attain-coupling} and
\(V_6(\eta)\le\gamma_\ell|\eta|-v_\ell(\eta!)\).
By Legendre's formula, the valuation of the multinomial coefficient
for the three parts is at most \(3f(\log b+2)\). Hence
\[
 v_\ell(l!)+V_{12}(l)\le
   \gamma_\ell |l|-\sum_\tau j_\tau+3f(\log b+2).
\]

In subtracting \(F(\ell)\), the difference between
\(\gamma_\ell\kappa(\tau)\) and the term of \(F(\ell)\) at \(\tau\) is
given by
\[
 v_\ell(\pi_\ell)\sum_{e\in I_\ell}\min\!\left(\frac{\kappa(\tau)}{\ell^e},\kappa(\tau_e)\right) .
\]
The chosen pair contributes at most \(j_\tau+1\).
For other \(2e<l_0\), one weight is \(a/x_1\), so the term is at most \(a/x_1\).
For \(2e\ge l_0\), summing the geometric series with
\(\lceil l_0/2\rceil-1\ge s\) and \(k\le2^s\) gives at most \(2b/(k\ell)\).
The errors per place after subtracting \(j_\tau\) are therefore bounded by
\[
 3(\log b+2)+x^{42}+1+l_0 a/x_1+2b/(k\ell)
 \le 3b/(k\ell)\le L/\ell-1\le v_\ell(B_\tau).
\]
For the first inequality the scales in Section~\ref{sec:data} give
\(3(\log b+2)+x^{42}+1\le a/x_1\) and
\((l_0+1)a/x_1\le b/(k\ell)\). The last counts multiples of \(\ell\)
in \(L\) consecutive integers. This proves \eqref{eq:attain-twelve}
at every prime.
\leansrc{New/Derivatives/Sharpness/Cases}{The factorial comparison in the three cases}.

\paragraph{Transfer to a coefficient of the theta kernel}
To extract the coefficient in \eqref{eq:attain-twelve}, the ultrametric
inequality lets us take each \(H_c\) to be a monomial in the \(y(\mu_i z)\)
of total degree at most \(|l|\). Each exponent \(e_{ci}\) satisfies
\(3e_{ci}\le3|l|\le6fb\le c_i\), so
\(y^{e_{ci}}\theta^{c_i}=(y\theta^3)^{e_{ci}}\theta^{c_i-3e_{ci}}\)
is a product of theta sections. The frame identity makes the twelvefold
product the theta series of a vector \(U\) in the Heisenberg model on
adelic \(12\times1\) matrices. Its finite values are integral, and it
is invariant under signed translations by \(\pi_M\mathfrak a_0^{12}\).
\leansrc{New/Derivatives/Sharpness/BoundAttained}{Transfer from theta sections to the Heisenberg model}.

Let \(I_{y_0}\) be the vacuum with indicator of the residue class \(y_0\)
modulo \((\pi_M)\) in its finite part.
Fourier inversion on the finite group expresses \(M^{12f}U\) as a sum
of translates of the vacuum. The coefficients \(a_j\) come from finite
sums of values of \(U\) weighted by roots of unity. More precisely,
for any finite value \(v_0\) of the vacuum, \(a_jv_0\) is integral.
Each translate is a sum of the \(I_{y_0}\) with roots of unity as
coefficients. For the linear map \(\Lambda\) taking a coefficient,
the ultrametric inequality thus gives some \(y_0\) for which
\(|\Lambda(I_{y_0})|_\iota\ge|M^{12f}v_0\Lambda(U)|_\iota\).
Include the integral vacuum normalization in \(\Lambda\), and fix a nonzero
value \(v_{\rm int}\) of the integral vacuum. By \eqref{eq:attain-twelve},
with
\[
 t=M^{12f}v_{\rm int}\prod_\tau B_\tau^2 ,
\]
some normalized coefficient of \(I_{y_0}\), multiplied by \(l!\), has
valuation at most \(F(\ell)+v_\ell(t)\).
Since \(B_\tau\le(2b+L)^L\) and all conjugates of \(v_{\rm int}\) have
modulus at most \(\exp(12C_{q_0,d}q_1^2)\), the bound
\(12f\log M+2fL\log(2b+L)+12C_{q_0,d}q_1^2\le fmb/k\) gives the
required bound on the conjugates of \(t\). The choice is independent
of \(\iota\).
\leansrc{New/Derivatives/Sharpness/BoundAttained}{Indicator coefficients and the bounds for conjugates}.
